\section{Discussion}
In this availability sample of 44 trials, human adjudication retained 413 of 440 candidates (93.9\%), including uncertain findings and protocol--report differences. The principal endpoint comprised 41 retained, confirmed or probable substantive statistical issues in report materials across 22 papers. Mean counts were 1.05 per paper in 2018--2019 and 0.84 in 2024--2025 (difference, $-0.21$; 95\% bootstrap confidence interval, $-1.02$ to $0.56$; exact permutation $P=0.654$). These data did not establish a period difference or equivalence.

TRACE-Q connects source locations, checking records, screening decisions, and candidate reports. These records allow reviewers to revisit a concern after the immediate working context has changed. This organization addresses the project's original practical motivation: checking source reports before their numbers enter subsequent research. It does not establish that the observed findings changed trial conclusions or downstream syntheses. Extraction artifacts must also be resolved against original documents before they become reporting concerns.

Candidate retention is not a measure of accuracy or sensitivity. Our adjudications concern findings already generated by the workflow, whereas planted-error studies such as that of La et al. assess recovery of specified targets.\cite{la2026consistency} Without an exhaustive reference set, a paper with no retained finding cannot be classified as free of problems. This application also formed part of workflow development, and exposure to earlier outputs may have influenced review rules. One reviewer supplied final adjudications; agreement between independent reviewers was not estimated. Judgments about context and substantive importance remain open to reassessment.

The period comparison describes the assembled article packages. Availability sampling and differences in clinical topics, trial designs, reporting complexity, and supplementary material may affect the opportunities for checking and the findings generated. Using papers as the units for permutation and bootstrap analyses preserves clustering of findings within papers, but does not remove selection effects or support causal claims about publication period, the pandemic, or AI use.

Statistical checks require contextual interpretation. Confidence intervals and $P$ values may differ because they use different procedures or adjustment sets. Crude ratios do not generally verify adjusted estimates, and documented amendments may explain departures from planned analyses. Display-precision thresholds and clinical rules of thumb identify questions for review; they do not define universal error boundaries. Dong et al.'s findings reinforce the importance of examining extraction and omitted analyses.\cite{dong2026statistical} Without participant-level data or analysis code, some concerns may remain unresolved.

We did not compare TRACE-Q with unaided human review, a single agent, or alternative agent environments, and did not establish time or cost savings. Agents using the same model may share mistaken assumptions. In a preprint, Kim et al. found that the effects of additional agents depended on task structure and coordination.\cite{kim2026scaling} Controlled comparisons with matched sources, models, tools, and computational budgets are needed, together with repeated runs to assess stability.

Workflow and analysis code can support reuse, but articles, supplements, and derived evidence may have redistribution restrictions. A reproducible release should identify source versions, workflow and model settings, the adjudication codebook, and source locators while distinguishing shareable materials from documents readers must obtain separately. Future evaluation should combine independently adjudicated natural findings with known inserted errors, report failed and incomplete reviews, and measure human verification time. TRACE-Q provides an exploratory framework for organizing quantitative review; its value in routine use requires prospective evaluation.
